\documentclass[../../../thesis.tex]{subfiles}
\begin{document}
\subsection{Tabuľky}
Sumarizácia dát vo forme tabuliek prispieva
k~sprehľadneniu obsahu, zjednodušuje text
a~umožňuje autorovi zamerať sa pri formulácii myšlienok na
obsahovú stránku práce.
Tabuľka, podobne ako obrázok, patrí medzi plávajúce objekty
a preto nemusí byť umiestnená priamo na mieste v dokumente,
kde sa o nej zmieňuje text.
Zvyčajne ju umiestňujeme za odsek s~prvou zmienkou,
ale často býva aj súčasťou prílohy dokumentu,
najmä ak je rozsiahlejšia.

Zameriame sa teraz iba na tabuľky s~výsledkami meraní,
ktoré sa v~záverečných prácach vyskytujú najčastejšie.

Tabuľku označujeme slovom Tabuľka,
za ktorým nasleduje poradové číslo tabuľky podľa výskytu
v~texte.
Za číslom môže nasledovať dvojbodka a~text s~opisom obsahu
tabuľky.
V~prípade, že označenie neobsahuje opisný text,
dvojbodku vynecháme.
Opisný text nekončí bodkou,
ani iným interpunkčným znamienkom,
pokiaľ ide iba o názov alebo jednu oznamovaciu vetu.
Celý odsek s~označením, číslom a~opisom umiestňujeme
nad tabuľku (pozri napríklad tabuľku~\ref{tab:template}).

\begin{table}[h!]
  \caption{Vzorová tabuľka}
  \label{tab:template}
  \centering
  \begin{tabular}{@{}lrrrr@{}}
    \toprule
    \textbf{názov riadka} & \textbf{stĺpec 1} & \textbf{stĺpec 2} & \textbf{stĺpec 3} & \textbf{stĺpec 4} \\
    \midrule
    prvý riadok & hodnota 1 & hodnota 2 & hodnota 3 & hodnota 4\\
    druhý riadok & hodnota 5 & hodnota 6 & hodnota 7 & hodnota 8\\
    tretí riadok & hodnota 9 & hodnota 10 & hodnota 11 & hodnota 12\\
    \bottomrule
  \end{tabular}
\end{table}

Na tabuľky sa odvolávame pomocou ich čísla použitím dvojice makier
\verb|\label| a~\verb|\ref|.
Spôsob odkazovania je podobný ako v prípade obrázkov,
o~ktorej sme podrobne hovorili v časti \ref{sec:figPlacement}.

\subsubsection{Vzhľad tabuľky}
Jednoduchá tabuľka obsahuje hlavičku a~niekoľko
údajových riadkov.
Vzhľad tabuľky je otázka estetických preferencií autora.
Príliš veľa grafických prvkov znižuje obsahovú hodnotu a čitateľnosť tabuľky.
Formát, ktorý sme vybrali je inšpirovaný trendmi
v~knižnej sadzbe.
Tabuľka je zhora a~zdola ohraničená vodorovnými čiarami
\verb|\toprule| a~\verb|\bottomrule| z~knižnice \verb|booktabs|.
Podobne je čiarou \verb|\midrule| oddelená hlavička tabuľky
a~prípadne aj päta, ak ju použijeme.
Zvislé čiary sa používajú iba vo výnimočných prípadoch,
napríklad ak je tabuľka rozdelená na
dve evidentne oddelené časti.
Prípadne môžeme čiarou oddeliť prvý stĺpec s~opisom označenia riadka (tabuľka \ref{tab:LED}).
Jednotlivé riadky s~údajmi neoddeľujeme.
Tabuľka pôsobí harmonicky,
ak je text v~prvom stĺpci zarovnaný doľava
a~v~poslednom stĺpci doprava.

Prvý riadok môže byť vysádzaný polotučným
rezom (bold),
ak obsahuje tzv. hlavičku, teda názvy stĺpcov alebo názvy
a~jednotky veličín, ktorých hodnoty sú v~konkrétnom stĺpci.
Označenie veličín symbolom ($U$, $I$, $R$, $P$, a~pod.)
nepíšeme v~hlavičke tučným písmom,
aby sme dodržali pravidlo o~tom,
že veličiny by mali byť v~celom dokumente označené symbolom
rovnakého tvaru a typu.

\begin{table}[h!]
  \caption{Tabuľka parametrov štyroch diód LED.
  FWHM predstavuje šírku píku v~polovici intenzity
  spektrálneho maxima (\foreignlanguage{english}{\emph{Full-Width-Half-Maximum}}) pri vlnovej dĺžke~$\lambda_\mathrm m$.}
  \label{tab:LED}
  \centering
  \begin{tabular}{@{}l|rrrr@{}}
    \toprule
    dióda & $\lambda_\mathrm m\,(\mathrm{nm})$ & FWHM\,(nm) & žiarivý výkon\,$(10^{-5}\,\mathrm W)$ & farba \\
    \midrule
    LED 1 & $450 \pm 5$ & $20\pm2$ & $3\pm1$ & modrá\\
    LED 2 & $525 \pm 5$ & $25\pm3$ & $50\pm4$ & zelená\\
    LED 3 & $615 \pm 5$ & $15\pm1$ & $2\pm1$ & oranžová\\
    LED 4 & $630 \pm 5$ & $20\pm2$ & $5\pm1$ & červená\\
    \bottomrule
  \end{tabular}
\end{table}

\subsubsection{Obsah tabuľky}
Tabuľka s nameranými hodnotami obsahuje v~prvom riadku označenie
veličín a~to buď slovom alebo symbolom.
Za veličinou nasleduje jednotka v~okrúhlej zátvorke.
Hranaté zátvorky na tento účel nepoužívame.
Bezrozmerné relatívne veličiny uvádzame s~jednotkou (a. u.).
Ide o zaužívanú formu v~medzinárodnej vedeckej komunite na pomenovanie tzv. príslušnej jednotky (angl. \foreignlanguage{english}{arbitrary unit}).
Takto označujeme aj osi grafov rôznych relatívnych veličín.

Pred jednotkou môže byť označenie násobku
a~dielu a~to ako v~symbolickej forme (kA, nm, MW),
tak aj vo forme dekadického exponentu ($10^3$, $10^{-9}$, $10^6$).
Vyhýbame sa zápisom v~tvare 1E-3 alebo 10-3,
pretože sú mätúce.
Text hlavičky vlnová dĺžka ($10^{-7}$\,m) znamená,
že hodnoty v~celom stĺpci predstavujú veličinu vlnová dĺžka
a~sú uvedené v~jednotkách $10^{-7}$ metra.

Neistoty a~odchýlky zapisujeme k~hlavnej hodnote pomocou
znaku $\pm$ alebo do zvláštneho stĺpca,
ktorý príslušne označíme.
Ďalšie spôsoby zápisu neistôt uvádza príslušná norma (napr. STN 01 6910: 2022 \cite{stn016910}).
\end{document}